\chapter{Discussion}
\label{ch:discussion}

\section{What the Instrument Test Shows}
\label{sec:discussion-test}
\highlight[orange!25]{The test asked whether the method holds together.}\todo{Based on AI-prepared stand-ins and an
invented walkthrough. Rewrite once the author has coded the policy and run a real session.} It does: a generic model,
built from the literature, was turned into trees; the trees were narrowed to one app by evidence from its own policy;
the text was coded against the same nodes; and each research question received an answer or a clear reason for having
none. Each step used the output of the one before it.

\paragraph{What the case suggests.} The aim of this thesis is to see which downstream harms providers acknowledge or
downplay. In this policy, the answer is that the chain breaks at its last link. Components and capabilities are stated,
and sometimes plainly, but the harm they could lead to is left for the reader to work out. A reader learns that a partner
receives an advertising identifier; a reader does not learn what that could mean. This is what the harm-tree lens was
meant to make visible, and a coding of data practices alone would not have shown it. It remains one policy, chosen on
purpose, read through stand-in coding. It suggests what to look for, and does not show that the pattern is common.

\paragraph{What the lens may overstate.} A privacy policy is a legal document, and naming harms may not be its role. The
gap may lie elsewhere: in help pages, in onboarding, or at the moment a user decides. This is why RQ4 matters, and why the
invented record cannot stand in for real sessions. Absence here means \enquote{not identified in the captured policy}. It
says nothing about intent or about what the provider does.

\paragraph{What has to change before the pilot.} The test found concrete weaknesses in the instrument:
\begin{itemize}[itemsep=1pt,topsep=3pt]
  \item \textbf{Harm category rule.} One named harm lifts every pathway of its category to level~3, even those it does not
    concern. The pilot should test tying a named harm to the capability the text links it to.
  \item \textbf{Node mapping.} The model maps mentions to nodes without regard to the data of the pathway. Reviewers need the
    pathway's data in view, and the prompt should state it.
  \item \textbf{Boundaries at stage~3.} The unstable framing and strength codes, and the many routine retention sentences coded
    as commitments, need sharper definitions and counter-examples.
  \item \textbf{Unsure flag.} At about a quarter of all labels, it would send too much to manual review to be useful.
  \item \textbf{Outside events.} Two of the four uncertain capabilities depend on events no material can settle. The record
    needs a fixed rule for what evidence would move them.
  \item \textbf{A real run.} The study needs a model with settable parameters, so that the settings can be reported as the method
    requires.
\end{itemize}

\section{Answering the Research Questions}
\todo[inline]{To be written.}

\section{Implications for Users}
\todo[inline]{To be written.}

\section{Implications for Providers and Designers}
\todo[inline]{To be written.}

\section{Implications for Regulation}
\todo[inline]{To be written.}

\section{Threats to Validity}
\todo[inline]{To be written.}
